\documentclass[a4paper]{article}
% generated by Docutils <http://docutils.sourceforge.net/>
\usepackage{fixltx2e} % LaTeX patches, \textsubscript
\usepackage{cmap} % fix search and cut-and-paste in Acrobat
\usepackage{ifthen}
\usepackage[T1]{fontenc}
\usepackage[utf8]{inputenc}
\usepackage{amsmath}
\usepackage[british,french,ngerman,english]{babel}
% Prevent side-effects if French hyphenation patterns are not loaded:
\frenchbsetup{StandardLayout}
\AtBeginDocument{\selectlanguage{english}\noextrasfrench}
\usepackage{color}
\usepackage{float} % float configuration
\floatplacement{figure}{H} % place figures here definitely
\usepackage{graphicx}
\usepackage{alltt}
\usepackage{multirow}
\usepackage{pifont}
\setcounter{secnumdepth}{0}
\usepackage{longtable,ltcaption,array}
\setlength{\extrarowheight}{2pt}
\newlength{\DUtablewidth} % internal use in tables
\usepackage{tabularx}
\usepackage{textcomp} % text symbol macros

%%% Custom LaTeX preamble
% PDF Standard Fonts
\usepackage{mathptmx} % Times
\usepackage[scaled=.90]{helvet}
\usepackage{courier}

%%% User specified packages and stylesheets

%%% Fallback definitions for Docutils-specific commands

% providelength (provide a length variable and set default, if it is new)
\providecommand*{\DUprovidelength}[2]{
  \ifthenelse{\isundefined{#1}}{\newlength{#1}\setlength{#1}{#2}}{}
}

% abstract title
\providecommand*{\DUtitleabstract}[1]{\centering\textbf{#1}}

% admonition (specially marked topic)
\providecommand{\DUadmonition}[2][class-arg]{%
  % try \DUadmonition#1{#2}:
  \ifcsname DUadmonition#1\endcsname%
    \csname DUadmonition#1\endcsname{#2}%
  \else
    \begin{center}
      \fbox{\parbox{0.9\linewidth}{#2}}
    \end{center}
  \fi
}

\makeatletter
\@namedef{DUrolealign-center}{\centering}
\makeatother


% dedication topic
\providecommand{\DUtopicdedication}[1]{\begin{center}#1\end{center}}

% docinfo (width of docinfo table)
\DUprovidelength{\DUdocinfowidth}{0.9\linewidth}

% subtitle (in document title)
\providecommand*{\DUdocumentsubtitle}[1]{{\large #1}}
\newcounter{enumv}

% error admonition title
\providecommand*{\DUtitleerror}[1]{\DUtitle{\color{red}#1}}

% fieldlist environment
\ifthenelse{\isundefined{\DUfieldlist}}{
  \newenvironment{DUfieldlist}%
    {\quote\description}
    {\enddescription\endquote}
}{}
% numeric or symbol footnotes with hyperlinks
\providecommand*{\DUfootnotemark}[3]{%
  \raisebox{1em}{\hypertarget{#1}{}}%
  \hyperlink{#2}{\textsuperscript{#3}}%
}
\providecommand{\DUfootnotetext}[4]{%
  \begingroup%
  \renewcommand{\thefootnote}{%
    \protect\raisebox{1em}{\protect\hypertarget{#1}{}}%
    \protect\hyperlink{#2}{#3}}%
  \footnotetext{#4}%
  \endgroup%
}

% inline markup (custom roles)
% \DUrole{#1}{#2} tries \DUrole#1{#2}
\providecommand*{\DUrole}[2]{%
  \ifcsname DUrole#1\endcsname%
    \csname DUrole#1\endcsname{#2}%
  \else% backwards compatibility: try \docutilsrole#1{#2}
    \ifcsname docutilsrole#1\endcsname%
      \csname docutilsrole#1\endcsname{#2}%
    \else%
      #2%
    \fi%
  \fi%
}

% legend environment
\ifthenelse{\isundefined{\DUlegend}}{
  \newenvironment{DUlegend}{\small}{}
}{}

% lineblock environment
\DUprovidelength{\DUlineblockindent}{2.5em}
\ifthenelse{\isundefined{\DUlineblock}}{
  \newenvironment{DUlineblock}[1]{%
    \list{}{\setlength{\partopsep}{\parskip}
            \addtolength{\partopsep}{\baselineskip}
            \setlength{\topsep}{0pt}
            \setlength{\itemsep}{0.15\baselineskip}
            \setlength{\parsep}{0pt}
            \setlength{\leftmargin}{#1}}
    \raggedright
  }
  {\endlist}
}{}

% optionlist environment
\providecommand*{\DUoptionlistlabel}[1]{\bf #1 \hfill}
\DUprovidelength{\DUoptionlistindent}{3cm}
\ifthenelse{\isundefined{\DUoptionlist}}{
  \newenvironment{DUoptionlist}{%
    \list{}{\setlength{\labelwidth}{\DUoptionlistindent}
            \setlength{\rightmargin}{1cm}
            \setlength{\leftmargin}{\rightmargin}
            \addtolength{\leftmargin}{\labelwidth}
            \addtolength{\leftmargin}{\labelsep}
            \renewcommand{\makelabel}{\DUoptionlistlabel}}
  }
  {\endlist}
}{}

% rubric (informal heading)
\providecommand*{\DUrubric}[2][class-arg]{%
  \subsubsection*{\centering\textit{\textmd{#2}}}}

% sidebar (text outside the main text flow)
\providecommand{\DUsidebar}[2][class-arg]{%
  \begin{center}
    \colorbox[gray]{0.80}{\parbox{0.9\linewidth}{#2}}
  \end{center}
}

% subtitle (for topic/sidebar)
\providecommand*{\DUsubtitle}[2][class-arg]{\par\emph{#2}\smallskip}

% title for topics, admonitions, unsupported section levels, and sidebar
\providecommand*{\DUtitle}[2][class-arg]{%
  % call \DUtitle#1{#2} if it exists:
  \ifcsname DUtitle#1\endcsname%
    \csname DUtitle#1\endcsname{#2}%
  \else
    \smallskip\noindent\textbf{#2}\smallskip%
  \fi
}

% titlereference role
\providecommand*{\DUroletitlereference}[1]{\textsl{#1}}

% topic (quote with heading)
\providecommand{\DUtopic}[2][class-arg]{%
  \ifcsname DUtopic#1\endcsname%
    \csname DUtopic#1\endcsname{#2}%
  \else
    \begin{quote}#2\end{quote}
  \fi
}

% transition (break, fancybreak, anonymous section)
\providecommand*{\DUtransition}[1][class-arg]{%
  \hspace*{\fill}\hrulefill\hspace*{\fill}
  \vskip 0.5\baselineskip
}

% hyperlinks:
\ifthenelse{\isundefined{\hypersetup}}{
  \usepackage[colorlinks=true,linkcolor=blue,urlcolor=blue]{hyperref}
  \urlstyle{same} % normal text font (alternatives: tt, rm, sf)
}{}
\hypersetup{
  pdftitle={reStructuredText Test Document},
  pdfauthor={David Goodger;Me;Myself;I}
}

%%% Title Data
\title{\phantomsection%
  reStructuredText Test Document%
  \label{restructuredtext-test-document}%
  \label{doctitle}%
  \\ % subtitle%
  \DUdocumentsubtitle{Examples of Syntax Constructs}%
  \label{examples-of-syntax-constructs}%
  \label{subtitle}}
\author{}
\date{}

%%% Body
\begin{document}
\maketitle

% Docinfo
\begin{center}
\begin{tabularx}{\DUdocinfowidth}{lX}
\textbf{Author}: &
	David Goodger \\
\textbf{Address}: &
	{\raggedright
123 Example Street\\
Example, EX  Canada\\
A1B 2C3 } \\
\textbf{Contact}: &
	\href{mailto:goodger@python.org}{goodger@python.org} \\
\textbf{Author}: &
	Me \\
\textbf{Author}: &
	Myself \\
\textbf{Author}: &
	I \\
\textbf{Organization}: &
	humankind \\
\textbf{Date}: &
	Now, or yesterday.  Or maybe even \emph{before} yesterday. \\
\textbf{Status}: &
	This is a “work in progress” \\
\textbf{Revision}: &
	is managed by a version control system. \\
\textbf{Version}: &
	1 \\
\textbf{Copyright}: &
	This document has been placed in the public domain. You
may do with it as you wish. You may copy, modify,
redistribute, reattribute, sell, buy, rent, lease,
destroy, or improve it, quote it at length, excerpt,
incorporate, collate, fold, staple, or mutilate it, or do
anything else to it that your or anyone else’s heart
desires. \\
\textbf{field name}: &
This is a “generic bibliographic field”.
\\
\textbf{field name “2”}: &
Generic bibliographic fields may contain multiple body elements.

Like this.
\\
\end{tabularx}
\end{center}

\DUtopic[dedication]{
\DUtitle[dedication]{Dedication}

For Docutils users \& co-developers.
}

\DUtopic[abstract]{
\DUtitle[abstract]{Abstract}

This is a test document, containing at least one example of each
reStructuredText construct.
}

% This is a comment. Note how any initial comments are moved by
% transforms to after the document title, subtitle, and docinfo.

% Above is the document title, and below is the subtitle.
% They are transformed from section titles after parsing.

% bibliographic fields (which also require a transform):

\pagebreak[4] % start ToC on new page

\phantomsection\label{table-of-contents}
\pdfbookmark[1]{Table of Contents}{table-of-contents}
\renewcommand{\contentsname}{Table of Contents}
\tableofcontents



\section{1~~~Structural Elements%
  \label{structural-elements}%
}


\subsection{1.1~~~Section Title%
  \label{section-title}%
}
\subsubsection*{Section Subtitle}

Lone subsections are converted to a section subtitle by a transform
activated with the \texttt{-{}-section-subtitles} command line option or the
\texttt{sectsubtitle-xform} configuration value.


\subsection{1.2~~~Empty Section%
  \label{empty-section}%
}


\subsection{1.3~~~Transitions%
  \label{transitions}%
}

Here’s a transition:


%___________________________________________________________________________
\DUtransition


It divides the section.  Transitions may also occur between sections:


%___________________________________________________________________________
\DUtransition



\section{2~~~Body Elements%
  \label{body-elements}%
}


\subsection{2.1~~~Paragraphs%
  \label{paragraphs}%
}

A paragraph.


\subsubsection{2.1.1~~~Inline Markup%
  \label{inline-markup}%
}

Paragraphs contain text and may contain inline markup: \emph{emphasis},
\textbf{strong emphasis}, \texttt{inline literals}, standalone hyperlinks
(\url{http://www.python.org}), external hyperlinks (\href{http://www.python.org/}{Python}\DUfootnotemark{id30}{id29}{5}), internal
cross-references (\hyperref[example]{example}), external hyperlinks with embedded URIs
(\href{http://www.python.org}{Python web site}), \href{http://www.python.org/}{anonymous hyperlink
references}\DUfootnotemark{id38}{id29}{5} (\href{http://docutils.sourceforge.net/}{a second reference}\DUfootnotemark{id40}{id39}{8}), footnote references (manually
numbered\DUfootnotemark{id1}{id8}{1}, anonymous auto-numbered\DUfootnotemark{id2}{id12}{3}, labeled auto-numbered\DUfootnotemark{id3}{label}{2}, or symbolic\DUfootnotemark{id4}{id13}{*}), citation references (\hyperlink{cit2002}{[CIT2002]}),
substitution references (\includegraphics{../../../docs/user/rst/images/biohazard.png}), and %
\phantomsection\label{inline-hyperlink-targets}inline hyperlink targets
(see \hyperref[targets]{Targets} below for a reference back to here).  Character-level
inline markup is also possible (although exceedingly ugly!) in \emph{re}\texttt{Structured}\emph{Text}.  Problems are indicated by %
\raisebox{1em}{\hypertarget{id28}{}}\hyperlink{id27}{\textbf{\color{red}|problematic|}} text
(generated by processing errors; this one is intentional).  Here is a
reference to the \hyperref[doctitle]{doctitle} and the \hyperref[subtitle]{subtitle}.

The default role for interpreted text is \DUroletitlereference{Title Reference}.  Here are
some explicit interpreted text roles: a PEP reference (\href{http://www.python.org/dev/peps/pep-0287}{PEP 287}); an
RFC reference (\href{http://tools.ietf.org/html/rfc2822.html}{RFC 2822}); an abbreviation (\DUrole{abbreviation}{abb.}), an acronym
(\DUrole{acronym}{reST}), code (\texttt{\DUrole{code}{print \textquotedbl{}hello world\textquotedbl{}}}); a \textsubscript{subscript};
a \textsuperscript{superscript} and explicit roles for \DUroletitlereference{Docutils}’
\emph{standard} \textbf{inline} \texttt{markup}.

% DO NOT RE-WRAP THE FOLLOWING PARAGRAPH!

Let’s test wrapping and whitespace significance in inline literals:
\texttt{This is an example of -{}-inline-literal -{}-text, -{}-including some-{}-
strangely-{}-hyphenated-words. ~Adjust-the-width-of-your-browser-window
to see how the text is wrapped. ~-{}- -{}-{}-{}- -{}-{}-{}-{}-{}-{}-{}- ~Now note ~ ~the
spacing ~ ~between the ~ ~words of ~ ~this sentence ~ ~(words
should ~ ~be grouped ~ ~in pairs).}

If the \texttt{-{}-pep-references} option was supplied, there should be a
live link to PEP 258 here.


\subsection{2.2~~~Bullet Lists%
  \label{bullet-lists}%
}
%
\begin{itemize}

\item A bullet list
%
\begin{itemize}

\item Nested bullet list.

\item Nested item 2.

\end{itemize}

\item Item 2.

Paragraph 2 of item 2.
%
\begin{itemize}

\item Nested bullet list.

\item Nested item 2.
%
\begin{itemize}

\item Third level.

\item Item 2.

\end{itemize}

\item Nested item 3.

\item This nested list should be compacted by the HTML writer.
%
\phantomsection\label{target}
% Even if this item contains a target and a comment.

\end{itemize}

\end{itemize}


\subsection{2.3~~~Enumerated Lists%
  \label{enumerated-lists}%
}
\begin{enumerate}

\item Arabic numerals.
\begin{enumerate}
\renewcommand{\labelenumii}{\alph{enumii})}

\item lower alpha)
\begin{enumerate}
\renewcommand{\labelenumiii}{(\roman{enumiii})}

\item (lower roman)
\begin{enumerate}

\item upper alpha.
\begin{list}{\Roman{enumv})}{\usecounter{enumv}}

\item upper roman)
\end{list}
\end{enumerate}
\end{enumerate}
\end{enumerate}

\item Lists that don’t start at 1:
\begin{enumerate}
\renewcommand{\labelenumii}{\arabic{enumii}.}
\setcounter{enumii}{2}

\item Three

\item Four
\end{enumerate}
\begin{enumerate}
\renewcommand{\labelenumii}{\Alph{enumii}.}
\setcounter{enumii}{2}

\item C

\item D
\end{enumerate}
\begin{enumerate}
\renewcommand{\labelenumii}{\roman{enumii}.}
\setcounter{enumii}{2}

\item iii

\item iv
\end{enumerate}
\end{enumerate}


\subsection{2.4~~~Definition Lists%
  \label{definition-lists}%
}
%
\begin{description}
\item[{Term}] \leavevmode 
Definition

\item[{Term}] \leavevmode (\textbf{classifier})

Definition paragraph 1.

Definition paragraph 2.

\item[{Term}] \leavevmode 
Definition

\item[{Term}] \leavevmode (\textbf{classifier one})
(\textbf{classifier two})

Definition

\end{description}


\subsection{2.5~~~Field Lists%
  \label{field-lists}%
}
%
\begin{DUfieldlist}
\item[{what:}]
Field lists map field names to field bodies, like database
records.  They are often part of an extension syntax.  They are
an unambiguous variant of RFC 2822 fields.

\item[{how arg1 arg2:}]
The field marker is a colon, the field name, and a colon.

The field body may contain one or more body elements, indented
relative to the field marker.

\item[{credits:}]
\DUrole{credits}{This paragraph has the \DUroletitlereference{credits} class set.  (This is actually not
about credits but just for ensuring that the class attribute
doesn’t get stripped away.)}

\end{DUfieldlist}


\subsection{2.6~~~Option Lists%
  \label{option-lists}%
}

For listing command-line options:
%
\begin{DUoptionlist}

\item[-a]  command-line option “a”

\item[-b file]  options can have arguments
and long descriptions

\item[-{}-long]  options can be long also

\item[-{}-input=file]  long options can also have
arguments

\item[-{}-very-long-option]  The description can also start on the next line.

The description may contain multiple body elements,
regardless of where it starts.

\item[-x, -y, -z]  Multiple options are an “option group”.

\item[-v, -{}-verbose]  Commonly-seen: short \& long options.

\item[-1 file, -{}-one=file, -{}-two file]  Multiple options with arguments.

\item[/V]  DOS/VMS-style options too

\end{DUoptionlist}

There must be at least two spaces between the option and the
description.


\subsection{2.7~~~Literal Blocks%
  \label{literal-blocks}%
}

Literal blocks are indicated with a double-colon (“::”) at the end of
the preceding paragraph (over there \texttt{-{}->}).  They can be indented:
%
\begin{quote}
\begin{alltt}
if literal_block:
    text = 'is left as-is'
    spaces_and_linebreaks = 'are preserved'
    markup_processing = None
\end{alltt}
\end{quote}

Or they can be quoted without indentation:
%
\begin{quote}
\begin{alltt}
>> Great idea!
>
> Why didn't I think of that?
\end{alltt}
\end{quote}


\subsection{2.8~~~Line Blocks%
  \label{line-blocks}%
}

This section tests line blocks.  Line blocks are body elements which
consist of lines and other line blocks.  Nested line blocks cause
indentation.

\begin{DUlineblock}{0em}
\item[] This is a line block.  It ends with a blank line.
\item[]
\begin{DUlineblock}{\DUlineblockindent}
\item[] New lines begin with a vertical bar (“|”).
\item[] Line breaks and initial indent are significant, and preserved.
\item[]
\begin{DUlineblock}{\DUlineblockindent}
\item[] Continuation lines are also possible.  A long line that is intended
to wrap should begin with a space in place of the vertical bar.
\end{DUlineblock}
\item[] The left edge of a continuation line need not be aligned with
the left edge of the text above it.
\end{DUlineblock}
\end{DUlineblock}

\begin{DUlineblock}{0em}
\item[] This is a second line block.
\item[] 
\item[] Blank lines are permitted internally, but they must begin with a “|”.
\end{DUlineblock}

Another line block, surrounded by paragraphs:

\begin{DUlineblock}{0em}
\item[] And it’s no good waiting by the window
\item[] It’s no good waiting for the sun
\item[] Please believe me, the things you dream of
\item[] They don’t fall in the lap of no-one
\end{DUlineblock}

Take it away, Eric the Orchestra Leader!
%
\begin{quote}

\begin{DUlineblock}{0em}
\item[] A one, two, a one two three four
\item[] 
\item[] Half a bee, philosophically,
\item[]
\begin{DUlineblock}{\DUlineblockindent}
\item[] must, \emph{ipso facto}, half not be.
\end{DUlineblock}
\item[] But half the bee has got to be,
\item[]
\begin{DUlineblock}{\DUlineblockindent}
\item[] \emph{vis a vis} its entity.  D’you see?
\item[] 
\end{DUlineblock}
\item[] But can a bee be said to be
\item[]
\begin{DUlineblock}{\DUlineblockindent}
\item[] or not to be an entire bee,
\item[]
\begin{DUlineblock}{\DUlineblockindent}
\item[] when half the bee is not a bee,
\item[]
\begin{DUlineblock}{\DUlineblockindent}
\item[] due to some ancient injury?
\item[] 
\end{DUlineblock}
\end{DUlineblock}
\end{DUlineblock}
\item[] Singing…
\end{DUlineblock}

\end{quote}

A line block, like the following poem by Christian Morgenstern, can
also be centre-aligned:

\begin{DUlineblock}{0em}
\foreignlanguage{ngerman}{\DUrole{align-center}{
\item[] \textbf{Die Trichter}
\item[] 
\item[] Zwei Trichter wandeln durch die Nacht.
\item[] Durch ihres Rumpfs verengten Schacht
\item[] fließt weißes Mondlicht
\item[] still und heiter
\item[] auf ~ ihren
\item[] Waldweg
\item[] u. s.
\item[] w.
\item[] 
}}
\end{DUlineblock}


\subsection{2.9~~~Block Quotes%
  \label{block-quotes}%
}

Block quotes consist of indented body elements:
%
\begin{quote}

My theory by A. Elk.  Brackets Miss, brackets.  This theory goes
as follows and begins now.  All brontosauruses are thin at one
end, much much thicker in the middle and then thin again at the
far end.  That is my theory, it is mine, and belongs to me and I
own it, and what it is too.
\nopagebreak

\raggedleft —Anne Elk (Miss)

\end{quote}

The language of a quote (like any other object) can be specified by
a class attribute:

% 
%
\begin{quote}
\foreignlanguage{french}{
ReStructuredText est un langage de balisage léger utilisé
notamment dans la documentation du langage Python.
}
\end{quote}


\subsection{2.10~~~Doctest Blocks%
  \label{doctest-blocks}%
}
%
\begin{quote}
\begin{alltt}
>>> print 'Python-specific usage examples; begun with ">>>"'
Python-specific usage examples; begun with ">>>"
>>> print '(cut and pasted from interactive Python sessions)'
(cut and pasted from interactive Python sessions)
\end{alltt}
\end{quote}


\subsection{2.11~~~Footnotes%
  \label{footnotes}%
}
%
\DUfootnotetext{id8}{id1}{1}{%
A footnote contains body elements, consistently indented by at
least 3 spaces.

This is the footnote’s second paragraph.
}
%
\DUfootnotetext{label}{id3}{2}{\phantomsection\label{label}%
Footnotes may be numbered, either manually (as in\DUfootnotemark{id9}{id8}{1}) or
automatically using a “\#”-prefixed label.  This footnote has a
label so it can be referred to from multiple places, both as a
footnote reference (\DUfootnotemark{id10}{label}{2}) and as a \hyperref[label]{hyperlink reference}.
}
%
\DUfootnotetext{id12}{id2}{3}{%
This footnote is numbered automatically and anonymously using a
label of “\#” only.

This is the second paragraph.

And this is the third paragraph.
}
%
\DUfootnotetext{id13}{id4}{*}{%
Footnotes may also use symbols, specified with a “*” label.
Here’s a reference to the next footnote:\DUfootnotemark{id14}{id15}{†}.
}
%
\DUfootnotetext{id15}{id14}{†}{%
This footnote shows the next symbol in the sequence.
}
%
\DUfootnotetext{id16}{id16}{4}{%
Here’s an unreferenced footnote, with a reference to a
nonexistent footnote:%
\raisebox{1em}{\hypertarget{id46}{}}%
\raisebox{1em}{\hypertarget{id17}{}}\hyperlink{id45}{\textbf{\color{red}{[}5{]}\_}}.
}


\subsection{2.12~~~Citations%
  \label{citations}%
}
\begin{figure}[b]\raisebox{1em}{\hypertarget{cit2002}{}}[CIT2002]
Citations are text-labeled footnotes. They may be
rendered separately and differently from footnotes.
\end{figure}

Here’s a reference to the above, \hyperlink{cit2002}{[CIT2002]}, and a %
\raisebox{1em}{\hypertarget{id48}{}}%
\raisebox{1em}{\hypertarget{id19}{}}\hyperlink{id47}{\textbf{\color{red}{[}nonexistent{]}\_}}
citation.


\subsection{2.13~~~Targets%
  \label{targets}%
  \label{another-target}%
}

\phantomsection\label{example}
This paragraph is pointed to by the explicit “example” target. A
reference can be found under \hyperref[inline-markup]{Inline Markup}, above. \hyperref[inline-hyperlink-targets]{Inline
hyperlink targets} are also possible.

Section headers are implicit targets, referred to by name. See
\hyperref[targets]{Targets}, which is a subsection of \hyperref[body-elements]{Body Elements}.

Explicit external targets are interpolated into references such as
“\href{http://www.python.org/}{Python}\DUfootnotemark{id31}{id29}{5}”.

Targets may be indirect and anonymous.  Thus \hyperref[targets]{this phrase} may also
refer to the \hyperref[targets]{Targets} section.

Here’s a %
\raisebox{1em}{\hypertarget{id50}{}}\hyperlink{id49}{\textbf{\color{red}`hyperlink reference without a target`\_}}, which generates an
error.


\subsubsection{2.13.1~~~Duplicate Target Names%
  \label{duplicate-target-names}%
}

Duplicate names in section headers or other implicit targets will
generate “info” (level-1) system messages.  Duplicate names in
explicit targets will generate “warning” (level-2) system messages.


\subsubsection{2.13.2~~~Duplicate Target Names%
  \label{id21}%
}

Since there are two “Duplicate Target Names” section headers, we
cannot uniquely refer to either of them by name.  If we try to (like
this: %
\raisebox{1em}{\hypertarget{id52}{}}\hyperlink{id51}{\textbf{\color{red}`Duplicate Target Names`\_}}), an error is generated.


\subsection{2.14~~~Directives%
  \label{directives}%
}

\phantomsection\label{contents}
These are just a sample of the many reStructuredText Directives.  For
others, please see
\url{http://docutils.sourceforge.net/docs/ref/rst/directives.html}.


\subsubsection{2.14.1~~~Document Parts%
  \label{document-parts}%
}

An example of the “contents” directive can be seen above this section
(a local, untitled table of \hyperref[contents]{contents}) and at the beginning of the
document (a document-wide \hyperref[table-of-contents]{table of contents}).


\subsubsection{2.14.2~~~Images and Figures%
  \label{images-and-figures}%
}

An image directive (also clickable – a hyperlink reference):

\hyperref[directives]{\includegraphics{../../../docs/user/rst/images/title.png}}

Image with multiple IDs:

\includegraphics{../../../docs/user/rst/images/title.png}
\phantomsection\label{image-target-3}\label{image-target-2}\label{image-target-1}

A centered image:

\noindent\makebox[\linewidth][c]{\includegraphics{../../../docs/user/rst/images/biohazard.png}}

A left-aligned image:

\noindent{\includegraphics{../../../docs/user/rst/images/biohazard.png}\hfill}

This paragraph might flow around the image.
The specific behavior depends upon the style sheet and
the browser or rendering software used.

A right-aligned image:

\noindent{\hfill\includegraphics{../../../docs/user/rst/images/biohazard.png}}

This paragraph might flow around the image.
The specific behavior depends upon the style sheet and
the browser or rendering software used.

For inline images see \hyperref[substitution-definitions]{Substitution Definitions}.

Image size:

An image 2 em wide:

\includegraphics[width=2em]{../../../docs/user/rst/images/biohazard.png}

An image 2 em wide and 15 pixel high:

\includegraphics[height=15px,width=2em]{../../../docs/user/rst/images/biohazard.png}

An image occupying 50\% of the line width:

\includegraphics[width=0.500\linewidth]{../../../docs/user/rst/images/title.png}

An image 2 cm high:

\includegraphics[height=2cm]{../../../docs/user/rst/images/biohazard.png}

A \emph{figure} is an image with a caption and/or a legend.  With page-based output
media, figures might float to a different position if this helps the page
layout.

\begin{figure}
\noindent\makebox[\linewidth][c]{\includegraphics[width=258bp]{../../../docs/user/rst/images/title.png}}
\caption{Plaintext markup syntax and parser system.}
\begin{DUlegend}
\setlength{\DUtablewidth}{\linewidth}
\begin{longtable*}[c]{|p{0.156\DUtablewidth}|p{0.563\DUtablewidth}|}
\hline

re
 & 
Revised, revisited, based on ‘re’ module.
 \\
\hline

Structured
 & 
Structure-enhanced text, structuredtext.
 \\
\hline

Text
 & 
Well it is, isn’t it?
 \\
\hline
\end{longtable*}

This paragraph is also part of the legend.
\end{DUlegend}
\end{figure}

A left-aligned figure:

\begin{figure} % align = "left"
\noindent\makebox[\linewidth][c]{\includegraphics[width=40px]{../../../docs/user/rst/images/biohazard.png}}
\caption{This is the caption.}
\begin{DUlegend}
This is the legend.

The legend may consist of several paragraphs.
\end{DUlegend}
\end{figure}

This paragraph might flow around the figure.

The specific behavior depends upon the style sheet and the browser or
rendering software used.

A centered figure:

\begin{figure}
\noindent\makebox[\linewidth][c]{\includegraphics[width=40px]{../../../docs/user/rst/images/biohazard.png}}
\caption{This is the caption.}
\begin{DUlegend}
This is the legend.

The legend may consist of several paragraphs.
\end{DUlegend}
\end{figure}

This paragraph might flow around the figure.

The specific behavior depends upon the style sheet and the browser or
rendering software used.

A right-aligned figure:

\begin{figure} % align = "right"
\noindent\makebox[\linewidth][c]{\includegraphics[width=40px]{../../../docs/user/rst/images/biohazard.png}}
\caption{This is the caption.}
\begin{DUlegend}
This is the legend.

The legend may consist of several paragraphs.
\end{DUlegend}
\end{figure}

This paragraph might flow around the figure. The specific behavior depends
upon the style sheet and the browser or rendering software used.


\subsubsection{2.14.3~~~Admonitions%
  \label{admonitions}%
}

\DUadmonition[attention]{
\DUtitle[attention]{Attention!}

Directives at large.
}

\DUadmonition[caution]{
\DUtitle[caution]{Caution!}

Don’t take any wooden nickels.
}

\DUadmonition[danger]{
\DUtitle[danger]{!DANGER!}

Mad scientist at work!
}

\DUadmonition[error]{
\DUtitle[error]{Error}

Does not compute.
}

\DUadmonition[hint]{
\DUtitle[hint]{Hint}

It’s bigger than a bread box.
}

\DUadmonition[important]{
\DUtitle[important]{Important}
%
\begin{itemize}

\item Wash behind your ears.

\item Clean up your room.

\item Call your mother.

\item Back up your data.

\end{itemize}
}

\DUadmonition[note]{
\DUtitle[note]{Note}

This is a note.
}

\DUadmonition[tip]{
\DUtitle[tip]{Tip}

15\% if the service is good.
}

\DUadmonition[warning]{
\DUtitle[warning]{Warning}

Strong prose may provoke extreme mental exertion.
Reader discretion is strongly advised.
}

\DUadmonition[admonition-and-by-the-way]{
\DUtitle[admonition-and-by-the-way]{And, by the way…}

You can make up your own admonition too.
}


\subsubsection{2.14.4~~~Topics, Sidebars, and Rubrics%
  \label{topics-sidebars-and-rubrics}%
}

\emph{Sidebars} are like miniature, parallel documents.

\DUsidebar{
\DUtitle[title]{Sidebar Title}

\DUsubtitle[sidebar]{Optional Subtitle}

This is a sidebar.  It is for text outside the flow of the main
text.

\DUrubric{This is a rubric inside a sidebar}

Sidebars often appear beside the main text with a border and a different
background or font color.
}

A \emph{topic} is like a block quote with a title, or a self-contained section
with no subsections.

\DUtopic[]{
\DUtitle[title]{Topic Title}

This is a topic.
}

A \emph{rubric} is like an informal heading that doesn’t correspond to the
document’s structure. It is typically highlighted in red (hence the name).

\DUrubric{This is a rubric}

Topics and rubrics can be used at places where a \hyperref[section-title]{section title} is not
allowed (e.g. inside a directive).


\subsubsection{2.14.5~~~Target Footnotes%
  \label{target-footnotes}%
}
%
\DUfootnotetext{id29}{id30}{5}{%
\url{http://www.python.org/}
}
%
\DUfootnotetext{id33}{id34}{6}{%
\url{http://pygments.org/}
}
%
\DUfootnotetext{id35}{id36}{7}{%
\url{ftp://ftp.ams.org/ams/doc/amsmath/short-math-guide.pdf}
}
%
\DUfootnotetext{id39}{id40}{8}{%
\url{http://docutils.sourceforge.net/}
}
%
\DUfootnotetext{id41}{id42}{9}{%
\url{A:DOS\\path\\}
}


\subsubsection{2.14.6~~~Replacement Text%
  \label{replacement-text}%
}

I recommend you try \href{http://www.python.org/}{Python, \emph{the} best language around}\DUfootnotemark{id32}{id29}{5}.


\subsubsection{2.14.7~~~Compound Paragraph%
  \label{compound-paragraph}%
}

Compound 1, paragraph 1.

Compound 1, paragraph 2.
%
\begin{itemize}

\item Compound 1, list item one.

\item Compound 1, list item two.

\end{itemize}

Another compound statement:

Compound 2, a literal block:
%
\begin{quote}
\begin{alltt}
Compound 2, literal.
\end{alltt}
\end{quote}
Compound 2, this is a test.

Compound 3, only consisting of one paragraph.
%
\begin{quote}
\begin{alltt}
Compound 4.
This one starts with a literal block.
\end{alltt}
\end{quote}
Compound 4, a paragraph.

Now something \emph{really} perverted – a nested compound block.  This is
just to test that it works at all; the results don’t have to be
meaningful.

Compound 5, block 1 (a paragraph).

Compound 6, block 2 in compound 5.

Compound 6, another paragraph.

Compound 5, block 3 (a paragraph).

Compound 7, with a table inside:

\setlength{\DUtablewidth}{\linewidth}
\begin{longtable*}[c]{|p{0.249\DUtablewidth}|p{0.249\DUtablewidth}|p{0.249\DUtablewidth}|}
\hline

Left cell, first
paragraph.

Left cell, second
paragraph.
 & 
Middle cell,
consisting of
exactly one
paragraph.
 & 
Right cell.

Paragraph 2.

Paragraph 3.
 \\
\hline
\end{longtable*}
Compound 7, a paragraph after the table.

Compound 7, another paragraph.


\subsubsection{2.14.8~~~Parsed Literal Blocks%
  \label{parsed-literal-blocks}%
}
%
\begin{quote}
\begin{alltt}
This is a parsed literal block.
    This line is indented.  The next line is blank.

Inline markup is supported, e.g. \emph{emphasis}, \textbf{strong}, \texttt{literal
text}, footnotes\DUfootnotemark{id22}{id8}{1}, %
\phantomsection\label{hyperlink-targets}hyperlink targets, and \href{http://www.python.org/}{references}.
\end{alltt}
\end{quote}


\subsubsection{2.14.9~~~Code%
  \label{code}%
}

Blocks of source code can be set with the \DUroletitlereference{code} directive. If the code
language is specified, the content is parsed and tagged by the \href{http://pygments.org/}{Pygments}\DUfootnotemark{id34}{id33}{6}
syntax highlighter and can be formatted with a style sheet. (Code parsing
is turned off using the \texttt{syntax-highlight} config setting in the test
conversions in order to get identical results with/without installed
Pygments highlighter.)
%
\begin{quote}
\begin{alltt}
print 'This is Python code.'
\end{alltt}
\end{quote}

The \texttt{:number-lines:} option (with optional start value) generates line
numbers:
%
\begin{quote}
\begin{alltt}
\DUrole{ln}{ 8 }# print integers from 0 to 9:
\DUrole{ln}{ 9 }for i in range(10):
\DUrole{ln}{10 }    print i
\end{alltt}
\end{quote}

For inline code snippets, there is the \DUroletitlereference{code} role, which can be used
directly (the code will not be parsed/tagged, as the language is not known)
or as base for special code roles, e.g. the LaTeX code in the next
paragraph.

Docutils uses LaTeX syntax for math directives and roles:
\texttt{\DUrole{code}{\DUrole{tex}{\textbackslash{}alpha = f(x)}}} prints $\alpha = f(x)$.

The \texttt{:code:} option of the \DUroletitlereference{include} directive sets the included content
as a code block, here the rst file \texttt{header\_footer.txt} with line numbers:
%
\begin{quote}
\begin{alltt}
\DUrole{ln}{1 }.. header:: Document header
\DUrole{ln}{2 }.. footer:: Document footer
\end{alltt}
\end{quote}


\subsection{2.15~~~Substitution Definitions%
  \label{substitution-definitions}%
}

An inline image (\includegraphics{../../../docs/user/rst/images/biohazard.png}) example:

(Substitution definitions are not visible in the HTML source.)


\subsection{2.16~~~Comments%
  \label{comments}%
}

Here’s one:

% Comments begin with two dots and a space. Anything may
% follow, except for the syntax of footnotes, hyperlink
% targets, directives, or substitution definitions.
% 
% Double-dashes -- "--" -- must be escaped somehow in HTML output.
% 
% Comments may contain non-ASCII characters: ä ö ü æ ø å

(View the HTML source to see the comment.)


\subsection{2.17~~~Raw text%
  \label{raw-text}%
}

This does not necessarily look nice, because there may be missing white space.

It’s just there to freeze the behavior.

A test.

Second test.

\DUrole{myclass}{Another test with myclass set.}

This is the \DUrole{myrawroleclass}{fourth test} with myrawroleclass set.

Fifth test in LaTeX.\\Line two.


\subsection{2.18~~~Container%
  \label{container}%
}

paragraph 1

paragraph 2

% currently not implemented in LaTeX:
% .. include:: data/header_footer.txt


\subsection{2.19~~~Colspanning tables%
  \label{colspanning-tables}%
}

This table has a cell spanning two columns:

\setlength{\DUtablewidth}{\linewidth}
\begin{longtable*}[c]{|p{0.075\DUtablewidth}|p{0.075\DUtablewidth}|p{0.086\DUtablewidth}|}
\hline
\multicolumn{2}{|p{0.15\DUtablewidth}|}{\textbf{%
Inputs
}} & \textbf{%
Output
} \\
\hline
\textbf{%
A
} & \textbf{%
B
} & \textbf{%
A or B
} \\
\hline
\endfirsthead
\hline
\multicolumn{2}{|p{0.15\DUtablewidth}|}{\textbf{%
Inputs
}} & \textbf{%
Output
} \\
\hline
\textbf{%
A
} & \textbf{%
B
} & \textbf{%
A or B
} \\
\hline
\endhead
\multicolumn{3}{c}{\hfill ... continued on next page} \\
\endfoot
\endlastfoot

False
 & 
False
 & 
False
 \\
\hline

True
 & 
False
 & 
True
 \\
\hline

False
 & 
True
 & 
True
 \\
\hline

True
 & 
True
 & 
True
 \\
\hline
\end{longtable*}


\subsection{2.20~~~Rowspanning tables%
  \label{rowspanning-tables}%
}

Here’s a table with cells spanning several rows:

\setlength{\DUtablewidth}{\linewidth}
\begin{longtable*}[c]{|p{0.296\DUtablewidth}|p{0.156\DUtablewidth}|p{0.226\DUtablewidth}|}
\hline
\textbf{%
Header row, column 1
(header rows optional)
} & \textbf{%
Header 2
} & \textbf{%
Header 3
} \\
\hline
\endfirsthead
\hline
\textbf{%
Header row, column 1
(header rows optional)
} & \textbf{%
Header 2
} & \textbf{%
Header 3
} \\
\hline
\endhead
\multicolumn{3}{c}{\hfill ... continued on next page} \\
\endfoot
\endlastfoot

body row 1, column 1
 & 
column 2
 & 
column 3
 \\
\hline

body row 2
 & \multirow{2}{0.16\DUtablewidth}{%
Cells may
span rows.
} & \multirow{2}{0.23\DUtablewidth}{%
Another
rowspanning
cell.
} \\
\cline{1-1}

body row 3
 &  \\
\hline
\end{longtable*}


\subsection{2.21~~~Custom Roles%
  \label{custom-roles}%
}
%
\begin{itemize}

\item A role based on an existing role.

\texttt{\DUrole{custom}{one}} \texttt{\DUrole{custom}{two}} \texttt{\DUrole{custom}{three}}

\item A new role.

\DUrole{customnew}{one two three}

\item A role with class attribute.

\DUrole{special}{interpreted text}

\item A language-switching role:

Let’s count in German \foreignlanguage{ngerman}{eins zwei drei}.

\item A role with multiple class attributes, styled with raw directives:

\newcommand{\DUrolegreen}[1]{\textcolor{green}{#1}}
\newcommand{\DUrolesc}[1]{\textsc{#1}}

The following works in most browsers but does not validate
(\texttt{<style>} is only allowed in the document head):
%
\begin{quote}
\begin{alltt}
.. raw:: html

  <style type="text/css"><!-{}-
   .green \{color: green;\}
   .sc \{font-variant: small-caps;\}
   -{}-></style>
\end{alltt}
\end{quote}

\DUrole{green}{\DUrole{sc}{\foreignlanguage{british}{British colourful text in small-caps}}}.

\end{itemize}


\subsection{2.22~~~Mathematics%
  \label{mathematics}%
}

Docutils supports inline math with the prefix or postfix \texttt{:math:}
role specificator, $n! + \sin(x_n^2)$ and $A_\text{c} =
\frac{\pi}{4} d^2$, as well as displayed math via the
\DUroletitlereference{math} directive:
%
\begin{equation*}
f(\epsilon) = \frac{1}{1 + \exp\left(\frac{\varepsilon}{k_\text{B}T}\right)}
\end{equation*}
Content may start on the first line of the directive, e.g.
%
\begin{equation*}
N = \frac{\text{number of apples}}{7}
\end{equation*}
Equations can be labeled with a reference name using the \texttt{:name:} option.
See \hyperref[eq-m]{eq:M} and \hyperref[eq-schrodinger]{eq:schrödinger} below.

The determinant of the matrix
%
\begin{equation*}
\mathbf{M} = \left(\begin{matrix}a&b\\c&d\end{matrix}\right)
\phantomsection
\label{eq-m}
\end{equation*}
is $|\mathbf{M}| = ad - bc$.

More than one display math block can be put in one math directive.
For example, the following sum and integral with limits:
%
\begin{equation*}
\int_0^1 x^n dx = \frac{1}{n + 1}
\end{equation*}%
\begin{equation*}
\sum_{n=1}^m n = \frac{m(m+1)}{2}
\end{equation*}
LaTeX-supported Unicode math symbols can be used in math roles and
directives:

The Schrödinger equation
%
\begin{equation*}
i\hbar \frac{\partial }{\partial t}\Psi  = \hat{H}\Psi ,
\phantomsection
\label{eq-schrodinger}
\end{equation*}
with the \emph{wave function} $\Psi $, describes how the quantum state of a
physical system changes in time.
%
\begin{description}
\item[{Math-Accents:}] \leavevmode 
\setlength{\DUtablewidth}{\linewidth}
\begin{longtable*}[c]{p{0.315\DUtablewidth}p{0.315\DUtablewidth}p{0.315\DUtablewidth}}

$\acute{a}$      \texttt{\textbackslash{}acute\{a\}}
 & 
$\dot{t}$        \texttt{\textbackslash{}dot\{t\}}
 & 
$\hat{\gamma}$   \texttt{\textbackslash{}hat\{\textbackslash{}gamma\}}
 \\

$\grave{a}$      \texttt{\textbackslash{}grave\{a\}}
 & 
$\ddot{t}$       \texttt{\textbackslash{}ddot\{t\}}
 & 
$\tilde{\alpha}$ \texttt{\textbackslash{}tilde\{\textbackslash{}alpha\}}
 \\

$\breve{x}$      \texttt{\textbackslash{}breve\{x\}}
 & 
$\dddot{t}$      \texttt{\textbackslash{}dddot\{t\}}
 & 
$\vec{\imath}$   \texttt{\textbackslash{}vec\{\textbackslash{}imath\}}
 \\

$\check{a}$      \texttt{\textbackslash{}check\{a\}}
 & 
$\bar{a}$        \texttt{\textbackslash{}bar\{a\}}
 & 
$\vec{R}$        \texttt{\textbackslash{}vec\{R\}}
 \\
\end{longtable*}

\end{description}

% \widetilde{xxx}
% \widehat{xxx}

Modulation Transfer Function:
%
\begin{equation*}
\text{MTF} = \left|\frac{\mathcal{F}\{s(x)\}}
                        {\mathcal{F}\{ s(x)\} |_{\omega _{x}=0}}\right|
           = \mathrm{abs}\left(\frac
                 {\int _{-\infty }^{\infty }s(x) \mathrm{e}^{\mathrm{i}\omega _{x}x}\mathrm{d}{x}}
                 {\int _{-\infty }^{\infty }s(x)\mathrm{d}{x}}
             \right).
\end{equation*}
Math split over two lines: If a double backslash is detected outside a
\texttt{\textbackslash{}begin\{...\} \textbackslash{}end\{...\}} pair, the math code is wrapped in an \href{ftp://ftp.ams.org/ams/doc/amsmath/short-math-guide.pdf}{AMSmath}\DUfootnotemark{id36}{id35}{7}
\texttt{align} environment:
%
\begin{align*}
s_{\mathrm{out}}(x) & = s_{\mathrm{in}}(x') * s_\delta (x-x') \\
                    & = \int  s_{\mathrm{in}}(x')s_\delta (x-x')\mathrm{d}x'
\end{align*}
Cases (“manually”, with \texttt{matrix} environment):
%
\begin{equation*}
\mathrm{sgn}(x) = \left\{\begin{matrix}
                    -1 & x<0\\
                    1 & x>0
                  \end{matrix}\right.
\end{equation*}
Cases with the \href{ftp://ftp.ams.org/ams/doc/amsmath/short-math-guide.pdf}{AMSmath}\DUfootnotemark{id37}{id35}{7} \texttt{cases} environment (not (yet) supported by
HTML writers with \texttt{-{}-math-output=MathML}):
%
\begin{equation*}
\mathrm{sgn}(x) = \begin{cases}
                    -1 & x<0\\
                     1 & x>0
                  \end{cases}
\end{equation*}

\section{3~~~Tests for the LaTeX writer%
  \label{tests-for-the-latex-writer}%
}

Test syntax elements which may cause trouble for the LaTeX writer but might
not need to be tested with other writers (e.g. the HTML writer).


\subsection{3.1~~~Custom Roles in LaTeX%
  \label{custom-roles-in-latex}%
}
%
\begin{itemize}

\item Role names and class arguments are converted to conform to the
regular expression \texttt{{[}a-z{]}{[}-a-z0-9{]}*} (letters are downcased,
accents and similar decoration is stripped, non-conforming
characters are replaced by a hyphen).

Class arguments may contain numbers and hyphens, which need special
treatment in LaTeX command names.

\DUrole{large}{\DUrole{custom4}{\DUrole{small-caps}{\DUrole{custom-role}{\DUrole{custom-role}{Interpreted Text}}}}}

\item With LaTeX, roles can be styled within the document using the \DUroletitlereference{raw}
directive.

\newcommand{\DUrolelarge}[1]{{\large #1}}
\makeatletter
\@namedef{DUrolesmall-caps}{\textsc}
\@namedef{DUrolecustom4}{\textbf}
\makeatother

\DUrole{large}{\DUrole{custom4}{\DUrole{small-caps}{\DUrole{custom-role}{\DUrole{custom-role}{Interpreted Text}}}}} in large, bold, small-caps.

\item Custom roles can be based on standard roles:

This is a \emph{\DUrole{custom-emphasis}{customized emphasis text role}}

This is a \texttt{\DUrole{custom-literal}{customized literal text role}}

This is a \textbf{\DUrole{custom-strong}{customized strong text role}}

This is a \textsubscript{\DUrole{custom-subscript}{customized subscript text role}}

This is a \textsuperscript{\DUrole{custom-superscript}{customized superscript text role}}

This is a \DUroletitlereference{\DUrole{custom-title-reference}{customized title-reference text role}}

\end{itemize}


\subsection{3.2~~~More Tables%
  \label{more-tables}%
}

A multicolumn table with multi-paragraph rowspanning cells:

\setlength{\DUtablewidth}{\linewidth}
\begin{longtable*}[c]{|p{0.133\DUtablewidth}|p{0.179\DUtablewidth}|p{0.179\DUtablewidth}|p{0.110\DUtablewidth}|p{0.121\DUtablewidth}|p{0.145\DUtablewidth}|}
\hline

test
 & 
\textbf{bold hd}
 & \multicolumn{3}{p{0.41\DUtablewidth}|}{
multicolumn 1

With a second paragraph
} & 
\emph{emph hd}
 \\
\hline
\multicolumn{2}{|p{0.31\DUtablewidth}|}{
multicolumn 2

With a second paragraph
} & 
cell
 & 
cell
 & 
cell
 & 
cell
 \\
\hline

cell
 & \multicolumn{2}{p{0.36\DUtablewidth}|}{
multicolumn 3 (one line,
but very very very very
very looooong)
} & 
cell
 & 
cell
 & 
cell
 \\
\hline

cell
 & 
cell
 & 
cell
 & \multicolumn{3}{p{0.38\DUtablewidth}|}{
Short multicolumn 4
} \\
\hline
\end{longtable*}

A table with multirow header
%
\begin{quote}

\setlength{\DUtablewidth}{\linewidth}
\begin{longtable*}[c]{|p{0.156\DUtablewidth}|p{0.238\DUtablewidth}|}
\hline
\multirow{2}{0.16\DUtablewidth}{%
\textbf{%
XXX
}} & \textbf{%
Variable Summary
} \\
\cline{2-2}
 & \textbf{%
Description
} \\
\hline
\endfirsthead
\hline
\multirow{2}{0.16\DUtablewidth}{%
\textbf{%
XXX
}} & \textbf{%
Variable Summary
} \\
\cline{2-2}
 & \textbf{%
Description
} \\
\hline
\endhead
\multicolumn{2}{c}{\hfill ... continued on next page} \\
\endfoot
\endlastfoot
\multicolumn{2}{|p{0.39\DUtablewidth}|}{
multirow header breaks latex
} \\
\hline
\end{longtable*}

\end{quote}

% This file is used by the standalone_rst_latex test.


\subsection{3.3~~~Option lists%
  \label{id23}%
}

The LaTeX-2e description environment is used for definition lists.
The definition is continued on the same line as the term, this should
not happen if a option-list is at the top of the definition.

If the option list is not at the first element in the definition, it
is contained in a quote
%
\begin{quote}
%
\begin{DUoptionlist}

\item[-{}-help]  show help

\item[-v]  verbose

\end{DUoptionlist}

\end{quote}
%
\begin{description}
\item[{In a definition list:}] \leavevmode %
\begin{DUoptionlist}

\item[-{}-help]  show help

\item[-v]  verbose

\end{DUoptionlist}

\end{description}


\subsection{3.4~~~Monospaced non-alphanumeric characters%
  \label{monospaced-non-alphanumeric-characters}%
}

These are all ASCII characters except a-zA-Z0-9 and space:

\texttt{!!!\textquotedbl{}\textquotedbl{}\textquotedbl{}\#\#\#\$\$\$\%\%\%\&\&\&'{}'{}'((()))***+++,{},{},-{}-{}-...///:::}

\texttt{;;;<{}<{}<===>{}>{}>???@@@{[}{[}{[}\textbackslash{}\textbackslash{}\textbackslash{}{]}{]}{]}\textasciicircum{}\textasciicircum{}\textasciicircum{}\_\_\_`{}`{}`\{\{\{|||\}\}\}\textasciitilde{}\textasciitilde{}\textasciitilde{}}

\texttt{xxxxxxxxxxxxxxxxxxxxxxxxxxxxxxxxxxxxxxxxxxxxxxxx}

The two lines of non-alphanumeric characters should both have the same
width as the third line.


\subsection{3.5~~~Non-ASCII characters%
  \label{non-ascii-characters}%
}

Punctuation and footnote symbols

\setlength{\DUtablewidth}{\linewidth}
\begin{longtable*}[c]{|p{0.028\DUtablewidth}|p{0.424\DUtablewidth}|}
\hline

–
 & 
en-dash
 \\
\hline

—
 & 
em-dash
 \\
\hline

‘
 & 
single turned comma quotation mark
 \\
\hline

’
 & 
single comma quotation mark
 \\
\hline

‚
 & 
low single comma quotation mark
 \\
\hline

“
 & 
double turned comma quotation mark
 \\
\hline

”
 & 
double comma quotation mark
 \\
\hline

„
 & 
low double comma quotation mark
 \\
\hline

†
 & 
dagger
 \\
\hline

‡
 & 
double dagger
 \\
\hline

\ding{169}
 & 
black diamond suit
 \\
\hline

\ding{170}
 & 
black heart suit
 \\
\hline

$\spadesuit$
 & 
black spade suit
 \\
\hline

$\clubsuit$
 & 
black club suit
 \\
\hline

…
 & 
ellipsis
 \\
\hline

™
 & 
trade mark sign
 \\
\hline

$\Leftrightarrow$
 & 
left-right double arrow
 \\
\hline
\end{longtable*}

The \DUroletitlereference{Latin-1 extended} Unicode block

\setlength{\DUtablewidth}{\linewidth}
\begin{longtable*}[c]{|p{0.051\DUtablewidth}|p{0.028\DUtablewidth}|p{0.028\DUtablewidth}|p{0.028\DUtablewidth}|p{0.028\DUtablewidth}|p{0.028\DUtablewidth}|p{0.028\DUtablewidth}|p{0.028\DUtablewidth}|p{0.028\DUtablewidth}|p{0.028\DUtablewidth}|p{0.028\DUtablewidth}|}
\hline

% 
 & 
0
 & 
1
 & 
2
 & 
3
 & 
4
 & 
5
 & 
6
 & 
7
 & 
8
 & 
9
 \\
\hline

160
 &  & 
¡
 & 
¢
 & 
£
 &  & 
¥
 & 
¦
 & 
§
 & 
¨
 & 
©
 \\
\hline

170
 & 
ª
 & 
«
 & 
¬
 & 
\-
 & 
®
 & 
¯
 & 
°
 & 
±
 & 
²
 & 
³
 \\
\hline

180
 & 
´
 & 
µ
 & 
¶
 & 
·
 & 
¸
 & 
¹
 & 
º
 & 
»
 & 
¼
 & 
½
 \\
\hline

190
 & 
¾
 & 
¿
 & 
À
 & 
Á
 & 
Â
 & 
Ã
 & 
Ä
 & 
Å
 & 
Æ
 & 
Ç
 \\
\hline

200
 & 
È
 & 
É
 & 
Ê
 & 
Ë
 & 
Ì
 & 
Í
 & 
Î
 & 
Ï
 & 
Ð
 & 
Ñ
 \\
\hline

210
 & 
Ò
 & 
Ó
 & 
Ô
 & 
Õ
 & 
Ö
 & 
×
 & 
Ø
 & 
Ù
 & 
Ú
 & 
Û
 \\
\hline

220
 & 
Ü
 & 
Ý
 & 
Þ
 & 
ß
 & 
à
 & 
á
 & 
â
 & 
ã
 & 
ä
 & 
å
 \\
\hline

230
 & 
æ
 & 
ç
 & 
è
 & 
é
 & 
ê
 & 
ë
 & 
ì
 & 
í
 & 
î
 & 
ï
 \\
\hline

240
 & 
ð
 & 
ñ
 & 
ò
 & 
ó
 & 
ô
 & 
õ
 & 
ö
 & 
÷
 & 
ø
 & 
ù
 \\
\hline

250
 & 
ú
 & 
û
 & 
ü
 & 
ý
 & 
þ
 & 
ÿ
 &  &  &  &  \\
\hline
\end{longtable*}
%
\begin{itemize}

\item The following line should not be wrapped, because it uses
no-break spaces (\textbackslash{}u00a0):

X~X~X~X~X~X~X~X~X~X~X~X~X~X~X~X~X~X~X~X~X~X~X~X~X~X~X~X~X~X~X~X~X~X~X~X~X~X~X~X~X~X~X~X~X~X~X~X~X~X~X~X~X~X~X~X~X~X~X~X~X~X~X~X~X~X

\item Line wrapping with/without breakpoints marked by soft hyphens
(\textbackslash{}u00ad):

pdn\-derd\-mdtd\-ri\-schpdn\-derd\-mdtd\-ri\-schpdn\-derd\-mdtd\-ri\-schpdn\-derd\-mdtd\-ri\-schpdn\-derd\-mdtd\-ri\-sch

pdnderdmdtdrischpdnderdmdtdrischpdnderdmdtdrischpdnderdmdtdrischpdnderdmdtdrisch

\item The currency sign (\textbackslash{}u00a4) is not supported by all fonts
(some have an Euro sign at its place). You might see an error
like:
%
\begin{quote}
\begin{alltt}
! Package textcomp Error: Symbol \textbackslash{}textcurrency not provided by
(textcomp)                font family ptm in TS1 encoding.
(textcomp)                Default family used instead.
\end{alltt}
\end{quote}

(which in case of font family ptm is a false positive). Add either
%
\begin{DUfieldlist}
\item[{warn:}]
turn the error in a warning, use the default symbol (bitmap), or

\item[{force,almostfull:}]
use the symbol provided by the font at the users
risk,

\end{DUfieldlist}

to the document options or use a different font package.

\end{itemize}


\subsection{3.6~~~Encoding special chars%
  \label{encoding-special-chars}%
}

The LaTeX Info pages lists under “2.18 Special Characters”
%
\begin{quote}

The following characters play a special role in LaTeX and are called
“special printing characters”, or simply “special characters”.
%
\begin{quote}

\# \$ \% \& \textasciitilde{} \_ \textasciicircum{} \textbackslash{} \{ \}

\end{quote}

\end{quote}

The special chars verbatim:
%
\begin{quote}
\begin{alltt}
# $ % & ~ _ ^ \textbackslash{} \{ \}
\end{alltt}
\end{quote}

However also \emph{square brackets} {[}{]} need special care.
%
\begin{quote}

item and all the other commands with optional arguments check
if the token right after the macro name is an opening bracket.
In that case the contents between that bracket and the following
closing bracket on the same grouping level are taken as the
optional argument. What makes this unintuitive is the fact that
the square brackets aren’t grouping characters themselves, so in
your last example item{[}{[}…{]}{]} the optional argument consists of
{[}… (without the closing bracket).

\end{quote}

Compare the items in the following lists:
%
\begin{itemize}

\item simple item

\item {[}bracketed{]} item

\end{itemize}
%
\begin{description}
\item[{simple}] \leavevmode 
description term

\item[{{[}bracketed{]}}] \leavevmode 
description term

\end{description}

The OT1 font-encoding has different characters for the less-than,
greater-than and bar, < | >, except for typewriter font \DUroletitlereference{cmtt}:
%
\begin{quote}
\begin{alltt}
< | >
\end{alltt}
\end{quote}


\subsection{3.7~~~Hyperlinks and -targets%
  \label{hyperlinks-and-targets}%
}

In LaTeX, we must set an explicit anchor (\texttt{\textbackslash{}phantomsection}) for a
%
\phantomsection\label{hypertarget-in-plain-text}hypertarget in plain text or in a figure but not in a longtable or
caption:

\setlength{\DUtablewidth}{\linewidth}
\begin{longtable}[c]{|p{0.075\DUtablewidth}|p{0.075\DUtablewidth}|p{0.075\DUtablewidth}|}
\caption{Table with %
\label{hypertarget-in-table-title}hypertarget in table title.}\\
\hline

False
 & 
True
 & 
None
 \\
\hline
\end{longtable}
\label{table-label}

\begin{figure}
\phantomsection\label{figure-label}
\noindent\makebox[\linewidth][c]{\includegraphics{../../../docs/user/rst/images/biohazard.png}}
\caption{Figure with %
\label{hypertarget-in-figure-caption}hypertarget in figure caption.}
\begin{DUlegend}
Legend with %
\phantomsection\label{hypertarget-in-figure-legend}hypertarget in figure legend.
\end{DUlegend}
\end{figure}

\includegraphics{../../../docs/user/rst/images/biohazard.png}
\phantomsection\label{image-label}

See \hyperref[hypertarget-in-plain-text]{hypertarget in plain text},
\hyperref[table-label]{table label}, \hyperref[hypertarget-in-table-title]{hypertarget in table title},
\hyperref[figure-label]{figure label}, \hyperref[hypertarget-in-figure-caption]{hypertarget in figure caption},
\hyperref[hypertarget-in-figure-legend]{hypertarget in figure legend}, and
\hyperref[image-label]{image label}.


\subsection{3.8~~~External references%
  \label{external-references}%
}

Long URLs should be wrapped in the PDF.
This can be achieved with the url command which is used by the LaTeX writer
whenever the content (name) of a reference node equals the link URL.
%
\begin{description}
\item[{Example:}] \leavevmode 
a long URL that should wrap in the output
\url{http://docutils.sourceforge.net/docs/user/latex.html\#id79}

\end{description}

If the argument contains any “\%”, “\#”, or “\textasciicircum{}\textasciicircum{}”, or ends with \texttt{\textbackslash{}}, it can’t
be used in the argument to another command. The argument must not contain
unbalanced braces.

The characters \textasciicircum{}, \{, \}, and \texttt{\textbackslash{}} are invalid in a “http:” or “ftp:” URL
and not recognized as part of it:

\begin{DUlineblock}{0em}
\item[] \url{http://www.example.org}/strange\textasciicircum{}\textasciicircum{}name
\item[] \url{http://www.example.org}\textbackslash{}using\textbackslash{}DOS\textbackslash{}paths\textbackslash{}
\item[] \url{http://www.example.org/XML}/strange\{n\}ame
\end{DUlineblock}

They can, however be used in paths and/or filenames.

Handling by the LaTeX writer:
%
\begin{itemize}

\item \texttt{\#}, \texttt{\textbackslash{}} and \texttt{\%} are escaped:

\begin{DUlineblock}{0em}
\item[] \href{http://www.w3.org/XML/Schema\#dev}{URL with \#}
\url{http://www.w3.org/XML/Schema\#dev}
\item[] \href{http://www.w3.org/XML/Schema\%dev}{URL with \%}
\url{http://example.org/Schema\%dev}
\item[] \href{A:DOS\\path\\}{file with DOS path}\DUfootnotemark{id42}{id41}{9} \url{A:DOS\\path\\}\DUfootnotemark{id43}{id41}{9}
\end{DUlineblock}

\DUadmonition[note]{
\DUtitle[note]{Note}

These URLs are typeset inside a LaTeX command without error.

\begin{DUlineblock}{0em}
\item[] \url{http://www.w3.org/XML/Schema\#dev}
\item[] \url{http://example.org/Schema\%dev}
\item[] \url{A:DOS\\path\\}\DUfootnotemark{id44}{id41}{9}
\end{DUlineblock}
}

\end{itemize}
%
\begin{itemize}

\item \textasciicircum{}\textasciicircum{} LaTeX’s special syntax for characters results in “strange” replacements
(both with href and url). A warning is given.

\href{../strange^^name}{file with \textasciicircum{}\textasciicircum{}}:
\url{../strange^^name}

\item Unbalanced braces, \{ or \}, will fail (both with href and url):
%
\begin{quote}
\begin{alltt}
`file with \{ <../strange\{name>`__
`<../strange\{name>`__
\end{alltt}
\end{quote}

while balanced braces are suported:

\begin{DUlineblock}{0em}
\item[] \url{../strange{n}ame}
\item[] \url{../st{r}ange{n}ame}
\item[] \url{../{st{r}ange{n}ame}}
\end{DUlineblock}

\end{itemize}


\subsection{3.9~~~Section titles with \hyperref[inline-markup]{inline markup}%
  \label{section-titles-with-inline-markup}%
}


\subsubsection{3.9.1~~~\emph{emphasized}, H\textsubscript{2}O and $x^2$%
  \label{emphasized-h2o-and-x-2}%
}


\subsubsection{3.9.2~~~Substitutions work%
  \label{substitutions-fail}%
}


\subsection{3.10~~~Deeply nested sections%
  \label{deeply-nested-sections}%
}

In LaTeX and HTML,


\subsubsection{3.10.1~~~Level 3%
  \label{level-3}%
}

nested sections


\paragraph{3.10.1.1~~~level 4%
  \label{level-4}%
}

reach at some level


\subparagraph{3.10.1.1.1~~~level 5%
  \label{level-5}%
}

(depending on the document class)


\DUtitle[sectionVI]{3.10.1.1.1.1~~~level 6%
  \label{level-6}%
}

an unsupported level.

% unusual combinations (from newlatex, for interactive testing)
% .. include:: data/latex.txt

% Preface for System Messages:


\section{4~~~Error Handling%
  \label{error-handling}%
}

Any errors caught during processing will generate system messages.

There should be five messages in the following, auto-generated
section, “Docutils System Messages”:

% section should be added by Docutils automatically


\section[Docutils System Messages]{\color{red}Docutils System Messages%
}

\DUadmonition[system-message]{
\DUtitle[system-message]{system-message}
\raisebox{1em}{\hypertarget{id27}{}}

{\color{red}ERROR/3} in \texttt{functional/input/data/standard.txt}, line~104

\hyperlink{id28}{
Undefined substitution referenced: \textquotedbl{}problematic\textquotedbl{}.
}}

\DUadmonition[system-message]{
\DUtitle[system-message]{system-message}
\raisebox{1em}{\hypertarget{id45}{}}

{\color{red}ERROR/3} in \texttt{functional/input/data/standard.txt}, line~391

\hyperlink{id46}{
Unknown target name: \textquotedbl{}5\textquotedbl{}.
}}

\DUadmonition[system-message]{
\DUtitle[system-message]{system-message}
\raisebox{1em}{\hypertarget{id47}{}}

{\color{red}ERROR/3} in \texttt{functional/input/data/standard.txt}, line~400

\hyperlink{id48}{
Unknown target name: \textquotedbl{}nonexistent\textquotedbl{}.
}}

\DUadmonition[system-message]{
\DUtitle[system-message]{system-message}
\raisebox{1em}{\hypertarget{id49}{}}

{\color{red}ERROR/3} in \texttt{functional/input/data/standard.txt}, line~427

\hyperlink{id50}{
Unknown target name: \textquotedbl{}hyperlink reference without a target\textquotedbl{}.
}}

\DUadmonition[system-message]{
\DUtitle[system-message]{system-message}
\raisebox{1em}{\hypertarget{id51}{}}

{\color{red}ERROR/3} in \texttt{functional/input/data/standard.txt}, line~440

\hyperlink{id52}{
Duplicate target name, cannot be used as a unique reference: \textquotedbl{}duplicate target names\textquotedbl{}.
}}

\end{document}
